% Ветка B, окончательная постановка (получена 2026-11-07). Дословная копия присланного текста.
% "2dimCDW"/"file 3dimCDW" = [CDW]; ссылки (130), (13), (12), Appendix A/D/G — на [CDW].
\documentclass[12pt,a4paper]{article}
\usepackage[utf8]{inputenc}
\usepackage[russian]{babel}
\usepackage{amsmath,amssymb,amsfonts,bm}
\usepackage{geometry}
\geometry{left=2cm,right=2cm,top=2cm,bottom=2cm}


\begin{document}

We consider a (2+1)-dimensional NJL-type model wih the following Lagrangian density,
\begin{equation}
L=\bar q\Big [\gamma^\rho\mathrm{i}D_\rho
+\mu\gamma^0\Big ]q+ \frac {G}{N}\Big
[(\bar qq)^2+(\bar q\mathrm{i}\gamma^5 q)^2 \Big ],
\label{1}
\end{equation}
where the fermion field $q(t,\vec r)$ (here $t\equiv x_0$ is a time-coordinate and $\vec r\equiv (x,y)$ is a radius-vector  of a point from the two-dimensional spatial plane)
as well as a reducible four-dimensional spinor multiplet of the (2+1)-dimensional Lorentz group, i.e. $q(t,\vec r)\equiv q_{si}(t,\vec r)$, where $s$ and $i$ are the spinor ($s=1,2,3,4$) and $U(N)$-group indices ($i=1,...,N$), correspondingly. Moreover, there the
summation over $\rho =0,1,2$ is implied, $\mu$ is the fermion number chemical potential. Expressions
as well as the algebra of the
$\gamma^\rho$ matrices ($\rho=0,1,2,5$) are the following.

{\bf $\gamma$ matrices.} The two-dimensional irreducible representation of the 3-dimensional
Lorentz group SO(2,1) is realized by the following $2\times 2$
$\tilde\gamma$-
matrices :
$$
\tilde\gamma^0=\sigma_3=
\left (\begin{array}{cc}
1 & 0\\
0 &-1
\end{array}\right ),\,\,
\tilde\gamma^1=i\sigma_1=
\left (\begin{array}{cc}
0 & i\\
i &0
\end{array}\right ),\,\,
\tilde\gamma^2=i\sigma_2=
\left (\begin{array}{cc}
0 & 1\\
-1 &0
\end{array}\right ),
\eqno (A.1)
$$
acting on two-component Dirac
spinors.

In addition to the  Dirac matrices $\gamma^\mu~~(\mu=0,1,2)$ there
exist other matrix, $\gamma^5$ which anticommute
with all $\gamma^\mu~~(\mu=0,1,2)$
\begin{eqnarray}
\gamma^5=\gamma^0\gamma^1\gamma^2\gamma^3=
i\left (\begin{array}{cc}
0~,& -I\\
I~,& 0
\end{array}\right ),\,\,
\label{A6}
\end{eqnarray}
with  $I$ being the unit $2\times 2$ matrix.




In our case the Dirac spinors $q$ have the following
form:
$$
q(x)=
\left (\begin{array}{cc}
\psi_{1}(x)\\
\psi_{2}(x)
\end{array}\right ),
\eqno (A.4)
$$
with $\psi_1,\psi_2$ being two-component spinors.
In the reducible four-dimensional spinor representation one deals
with
(4$\times$4) $\gamma$-matrices:
$\gamma^\mu=diag(\tilde\gamma^\mu,-\tilde\gamma^\mu)$, where
$\tilde\gamma^\mu$

where as well
$$
D_{\mu}=\partial_{\mu}+ieA_{\mu}
$$
 where $A_{\mu}=(A_0,A_1,A_2)$, in our case $A_0=A_1=0$ and $A_2=Bx_1$, $B=const$, we can fix the sign of $B$ or $eB$, where $x_1$ is the first coordinate from plane ($x_1$, $x_2$).

\section{ The model and its thermodynamic potential}
%\label{effaction}

For simplicity, first we consider a (2+1)-dimensional NJL-type model at $B=0$. It has the following Lagrangian density,
\begin{equation}
L=\bar q\Big [\gamma^\rho\mathrm{i}\partial_\rho
+\mu\gamma^0\Big ]q+ \frac {G}{N}\Big
[(\bar qq)^2+(\bar q\mathrm{i}\gamma^5 q)^2 \Big ],
\label{1}
\end{equation}
where the fermion field $q(t,\vec r)$ (here $t\equiv x^0$ is a time- and
$\vec r\equiv (x,y)$ is a vector of spatial coordinates of an arbitrary point from the
(2+1)-dimensional space-time) is transformed as a fundamental multiplet of the $O(N)$ group
\footnote{In the present consideration this group is a rather auxiliary character. It is introduced with the aim to study the model in terms of the leading order of the nonperturbative large $N$-method.}
as well as a reducible four-dimensional spinor multiplet of the (2+1)-dimensional Lorentz group,
i.e. $q(t,\vec r)\equiv q_{si}(t,\vec r)$, where $s$ and $i$ are the spinor ($s=1,2,3,4$) and
$O(N)$-group indices ($i=1,...,N$), respectively. Moreover, there the
summation over $\rho =0,1,2$ is implied, $\mu$ is the fermion number chemical potential, and the
corresponding expressions for Dirac spinors as well as the algebra of the
$\gamma^\rho$ matrices ($\rho=0,1,2,5$) are presented in Appendix A.  Clearly, the Lagrangian
(\ref{1}) is invariant under the action both of the above mentioned auxiliary $O(N)$ group and
transformations from the $U_V(1)\times U_A(1)$ group, where $U_V(1)$ is responsible for
the electric charge conservation, $q\to e^{i\alpha}q$, and $U_A(1)$ is the group of continuous
axial (chiral) transformations, $q\to e^{i\alpha\gamma^5}q$.

The linearized version of the Lagrangian (\ref{1}), which contains two auxiliary
composite bosonic fields (the scalar $\sigma (t,\vec r)$ and the pseudoscalar $\pi (t,\vec r)$),
has the following form:
\begin{eqnarray}
\widetilde L\ds &=&\bar q\Big [\gamma^\rho\mathrm{i}\partial_\rho
+\mu\gamma^0-\sigma -\mathrm{i}\gamma^5\pi\Big ]q
 -\frac{N}{4G}\Big [\sigma^2+\pi^2\Big ].
\label{2}
\end{eqnarray}
Evidently, the Lagrangian (\ref{2}) is equivalent to (\ref{1}), which simply follows from the use of the
following constraint equations for the bosonic fields
\begin{eqnarray}
\sigma(t,\vec r)=-2\frac G{N}(\bar qq);~~~\pi (t,\vec r)=-2\frac G{N}(\bar q
\mathrm{i}\gamma^5 q). \label{200}
\end{eqnarray}
Furthermore, it is clear from (\ref{200}) that the bosonic fields transform under the axial $U_{A}(1)$ group in the following manner:
\begin{eqnarray}
U_{A}(1):~~\sigma \to\cos (2\alpha)\sigma +\sin (2\alpha)\pi;
~~\pi\to\cos (2\alpha)\pi-\sin (2\alpha)\sigma.
\label{201}
\end{eqnarray}
Starting from Lagrangian (\ref{2}), one obtains in the leading order
of the large $N$-expansion (i.e. in the one-fermion loop
approximation) the following path integral expression for the
effective action ${\cal S}_{\rm {eff}}(\sigma,\pi)$ of the bosonic
$\sigma (t,\vec r)$ and $\pi (t,\vec r)$ fields:
$$
\exp(\mathrm{i}{\cal S}_{\rm {eff}}(\sigma,\pi))=
  N'\int[d\bar q][dq]\exp\Bigl(\mathrm{i}\int\widetilde L\,dtd^2\vec r\Bigr),
$$
where
\begin{equation}
{\cal S}_{\rm {eff}}
(\sigma,\pi)
=-N\int dtd^2\vec r\left [\frac{\sigma^2+\pi^2}{4G}
\right ]+\widetilde {\cal S}_{\rm {eff}},
\label{3}
\end{equation}
and $N'$ is a normalization constant.
The fermion contribution to the effective action, i.e. the term
$\widetilde {\cal S}_{\rm {eff}}$ in (\ref{3}), is given by:
\begin{equation}
\exp(\mathrm{i}\widetilde {\cal S}_{\rm {eff}})=N'\int [d\bar
q][dq]\exp\Bigl(\mathrm{i}\int\Big\{\bar q\big
[\gamma^\rho\mathrm{i}\partial_\rho +\mu\gamma^0-\sigma -\mathrm{i}\gamma^5\pi\big
]q\Big\}dtd^2\vec r\Bigr).
 \label{4}
\end{equation}
The ground state expectation values  $\langle\sigma(t,\vec r)\rangle$ and
$\langle\pi(t,\vec r)\rangle$, of the composite bosonic fields are determined by
the saddle point equations,
\begin{eqnarray}
\frac{\delta {\cal S}_{\rm {eff}}}{\delta\sigma (t,\vec r)}=0,~~~~~
\frac{\delta {\cal S}_{\rm {eff}}}{\delta\pi (t,\vec r)}=0.~~~~~
\label{5}
\end{eqnarray}

In particular, in this paper we will
use the following ansatz, which is usually called the chiral density wave (CDW) \cite{nakano} or chiral spiral \cite{thies} one:
\begin{eqnarray}
\langle\sigma(t,\vec r)\rangle=M,~~~\langle\pi(t,\vec r)\rangle=0, \label{6}
\end{eqnarray}
where $M$ is time- and $\vec r$- independent quantities.


In the leading order of the large $N$-expansion $\Omega_\mu (M,b)$ is defined by the following expression:
\begin{equation}
\int dtd^2\vec r \Omega_\mu (M)=-\frac{1}{N}{\cal S}_{\rm {eff}}\{\sigma(t,\vec r),\pi(t,\vec r)\}\big|_{\sigma (t,\vec r)=\vev{\sigma(t,\vec r)},\pi(t,\vec r)=\vev{\pi(t,\vec r)}} ,
\end{equation}
which gives
\begin{equation}
\mathrm{i}\int dtd^2\vec r\Omega_\mu (M,b)\,\,=\,\,\mathrm{i}\int
dtd^2\vec r\frac{M^2}{4G}-\frac{1}{N}\ln\left (
\int [d\bar q][dq]\exp\Bigl(\mathrm{i}\int dtd^2\vec r
\bar q {\cal D} q \Bigr)\right ),
\label{8}
\end{equation}
where
\begin{equation}
{\cal D}=\gamma^\rho\mathrm{i}\partial_\rho +\mu\gamma^0+
M.\label{9}
\end{equation}

First, let us study the properties of the TDP (\ref{130}) in the supposition that chiral condensates (\ref{6}) are spatially homogenious, i.e. there we have $b=0$. And in this case we have from Eq. (\ref{13}) the following expression for TDP
\begin{eqnarray}
\Omega_\mu(M)&
\equiv & \Omega^{un}_\mu(M)=\frac{M^2}{4G}-2\int\frac{d^2p}{
(2\pi)^2}\Big\{E+(\mu-E)\theta (\mu-E)\Big\},
\label{0134}
\end{eqnarray}
where $E=\sqrt{p_1^2+p_2^2+M^2}$. %\vspace{0.3cm}

\subsection{The case with $\mu=0$}
\label{IIIA}
%$\bullet$ {\bf The case with $\mu=0$.}

First of all let us consider the vacuum case, i.e. when
$\mu=0$ and temperature is zero.  Assuming homogeneity of
the chiral condensate ($b=0$), we obtain from (\ref{0134}) the
following expression for the unrenormalized effective potential of the initial
NJL$_3$ model in vacuum (note that in vacuum the TDP is usually called effective potential):
 \begin{eqnarray}
\Omega_{\mu=0}^{un}(M)=V^{un}(M)&=&\frac{M^2}{4G}-
\frac{2}{(2\pi)^2}\int d^2p\sqrt{p_1^2+p_2^2+M^2}.
\label{15}
\end{eqnarray}
Formally, the effective potential (\ref{15}) is an UV
divergent quantity, but recall that the NJL$_3$ model (1) is renormalizable \cite{rosenstein}. To renormalize the effective potential
$V^{un}(M)$, i.e. to obtain a finite expression $V^{ren}(M)$ for it, we first
need to regularize the integral in the right hand side of (\ref{15})
by cutting off its integration region. Below, we consider two
different regularization schemes.

$\bullet$
{\bf The momentum cutoff regularization}. First of all, let us regularize the expression (\ref{15}) by cutting
momenta, i.e. we suppose that $|p_1|<\Lambda$, $|p_2|<\Lambda$ in
(\ref{15}). As a result we have the following regularized expression
for the effective potential (which is finite at finite values of
$\Lambda$):
 \begin{eqnarray}
V^{reg1}(M)&=&\frac{M^2}{4G}-\frac{2}{\pi^2}\int_0^\Lambda
dp_1\int_0^\Lambda dp_2\sqrt{p_1^2+p_2^2+M^2} \nonumber\\&=&M^2\left
[\frac 1{4G}-\frac{2\Lambda\ln(1+\sqrt{2})}{\pi^2}\right
]+\frac{M^3}{3\pi}-
\frac{2\Lambda^3(\sqrt{2}+\ln(1+\sqrt{2}))}{3\pi^2}+{\cal
O}(1/\Lambda), \label{14}
\end{eqnarray}
where ${\cal O}(1/\Lambda)$ denotes an expression which tends to
zero as $1/\Lambda$ at $\Lambda\to \infty$.
 Second, we suppose that the bare coupling constant
$G$ in (\ref{15})-(\ref{14}) depends on the cutoff parameter $\Lambda$
%($G\equiv G(\Lambda)$)
in such a way that in the limit
$\Lambda\to\infty$ one obtains a finite expression for effective potential. To construct the
function $G$ vs $\Lambda$, let us look at the expression in the square brackets of (\ref{14}).
Then it is clear that
the dependence of $G$ on the cutoff $\Lambda$ should be the following one
 \begin{eqnarray}
\frac 1{4G}\equiv \frac
1{4G(\Lambda)}=\frac{2\Lambda\ln(1+\sqrt{2})}{\pi^2}+\frac{1}{2g},
\label{16}
\end{eqnarray}
where $g$ is a finite and $\Lambda$ independent model parameter with dimensionality of inverse mass. Hence, taking
into account in (\ref{14}) the relation (\ref{16}) and ignoring there an infinite $M$-independent constant,
one obtains the following renormalized expression for the effective potential $V(M)$, i.e.
\begin{eqnarray}
V^{ren}(M)&=&\lim_{\Lambda\to\infty}\left\{V^{reg1}
(M)\Big |_{G= G(\Lambda)}+\frac{2\Lambda^3(\sqrt{2}+\ln(1+\sqrt{2}))}{3\pi^2}\right\}%\nonumber\\
%&=&
=\frac{M^2}{2g}+\frac{M^3}{3\pi}. \label{17}
\end{eqnarray}
The fact that it is possible to renormalize the effective potential
of the initial model (1) in the leading order of the large-$N$ expansion is the reflection of a more general property of
(2+1)-dimensional theories with four fermion interactions. Indeed,
it is well known that in the framework of  the ''naive''
perturbation theory (over coupling constants) these models are not
renormalizable ones. However, as it was proved in \cite{rosenstein},
in the framework of nonperturbative large-$N$ technique these models
are renormalizable in each order of $1/N$ expansion.

It is clear from (\ref{17}) that at $g>0$ the global minimum of the
effective potential lies at $M=0$. So in this case chiral symmetry
remains intact and fermions are massless. However, if $g<0$ then the
global minimum of the function (\ref{17}) corresponds to the point
$M_0=-\pi/g$. In this case chiral symmetry is broken spontaneously
and fermions acquire dynamically the mass $M_0$.


\subsection{The case with $\mu\ne 0$}

It is easily seen from (\ref{13}) that at $\mu>0$ and $b=0$ the TDP of  the system looks like
\begin{eqnarray}
\Omega_\mu (M,b=0)&
=&\frac{M^2}{4G}-\frac{2}{(2\pi)^2}\int d^2p\sqrt{p_1^2+p_2^2+M^2}\nonumber\\
&-&\frac{2}{(2\pi)^2}\int d^2p\left
(\mu-\sqrt{p_1^2+p_2^2+M^2}\right ) \theta\left
(\mu-\sqrt{p_1^2+p_2^2+M^2}\right ). \label{20}
\end{eqnarray}
Evidently, after renormalization the first two terms in (\ref{20}) bring us to the
vacuum effective potential $V^{ren}(M)$ (\ref{17})-(\ref{19}), whereas in the
last integral of Eq. (\ref{20}), which is an UV-convergent one, it is easy
to integrate over $p_1$ and $p_2$. As a result, we obtain the following expression for the renormalized
TDP at $\mu>0$ and $b=0$ (see, e.g.,  in \cite{klimenko}):
\begin{eqnarray}
\Omega^{ren}_\mu (M,b=0)&=&V^{ren}(M)-\frac{1}{6\pi}
\theta (\mu-M)(\mu-M)^2(\mu+2M).
\label{21}\end{eqnarray}
A study of this TDP shows (see, e.g., in \cite{klimenko}) that at $g>0$, the chiral symmetry
$U_A(1)$ of the model (1) is not violated for any values of the chemical potential. However, at
$g<0$, in the model with $\mu=0$ there is a spontaneous generation of fermion
mass equal to $M_0=-\pi /g$. Then at $0\le\mu\le M_0$ the chiral symmetry is broken
spontaneously and fermions have nonzero mass $M_0$, but at $\mu=M_0$ a first order phase
transition occurs, and at $\mu >M_0$ a phase with massless fermions is realized in the model.



\section{Stability of homogeneous vacuum in 3D NJL-type model under small inhomogeneous
fluctuations}

For simplicity, we consider a (2+1)-dimensional NJL-type model (below, for it we use also the
notation NJL$_3$) with only one $O(N)$ multiplet of fermion fields. It has the following Lagrangian density,
\begin{equation}
L=\bar q\Big [\gamma^\rho\mathrm{i}\partial_\rho
+\mu\gamma^0\Big ]q+ \frac {G}{N}\Big
[(\bar qq)^2+(\bar q\mathrm{i}\gamma^5 q)^2 \Big ],
\label{F1}
\end{equation}
where the fermion field $q(t,\vec r)$ (here $t\equiv x^0$ is a time- and
$\vec r\equiv (x,y)$ is a vector of spatial coordinates of an arbitrary point from the
(2+1)-dimensional space-time) is transformed as a fundamental multiplet of the $O(N)$ group
\footnote{In the present consideration this group is a rather auxiliary character. It is
introduced with the aim to study the model in terms of the leading order of the nonperturbative
large $N$-method.}
as well as a reducible four-dimensional spinor multiplet of the (2+1)-dimensional Lorentz group,
i.e. $q(t,\vec r)\equiv q_{si}(t,\vec r)$, where $s$ and $i$ are the spinor ($s=1,2,3,4$) and
$O(N)$-group indices ($i=1,...,N$), respectively. Moreover, there the
summation over $\rho =0,1,2$ is implied, $\mu$ is the fermion number chemical potential. The
corresponding expressions for Dirac spinors as well as the algebra of the
$\gamma^\rho$ matrices ($\rho=0,1,2,5$) are presented in Appendix A.  Clearly, the Lagrangian
(\ref{F1}) is invariant under the action both of the above mentioned auxiliary $O(N)$ group and
transformations from the $U_V(1)\times U_A(1)$ group, where $U_V(1)$ is responsible for
the electric charge conservation, $q\to e^{i\alpha}q$, and $U_A(1)$ is the group of continuous
axial (chiral) transformations, $q\to e^{i\alpha\gamma^5}q$.

The linearized version of the Lagrangian (\ref{F1}), which contains two auxiliary
composite bosonic fields (the scalar $\sigma (t,\vec r)$ and the pseudoscalar $\pi_0 (t,\vec r)$),
has the following form:
\begin{eqnarray}
\widetilde L\ds &=&\bar q\Big [\gamma^\rho\mathrm{i}\partial_\rho
+\mu\gamma^0+\sigma +\mathrm{i}\gamma^5\pi_0\Big ]q
 -\frac{N}{4G}\Big [\sigma^2+\pi_0^2\Big ].
\label{F2}
\end{eqnarray}
Evidently, the Lagrangian (\ref{F2}) is equivalent to (\ref{F1}), which simply follows from the use of the
following constraint equations for the bosonic fields
\begin{eqnarray}
\sigma(t,\vec r)=2\frac G{N}(\bar qq);~~~\pi_0 (t,\vec r)=2\frac G{N}(\bar q
\mathrm{i}\gamma^5 q). \label{F200}
\end{eqnarray}
Moreover, it follows from Eqs. (\ref{F200}) that under continuous chiral transformations, when $q\to e^{i\alpha\gamma^5}q$, auxiliary bosonic fields are transformed as
\begin{eqnarray}
U_{A}(1):~~\sigma \to\cos (2\alpha)\sigma +\sin (2\alpha)\pi;
~~\pi\to\cos (2\alpha)\pi-\sin (2\alpha)\sigma.
\label{F201}
\end{eqnarray}
As a result, we see from Eq. (\ref{F201}) that the field combination $(\sigma^2(t,\vec r)+\pi_0^2 (t,\vec r))$ remains intact under chiral $U_A(1)$ transformations.
Starting from Lagrangian (\ref{F2}), one obtains in the leading order
of the large $N$-expansion (i.e. in the one-fermion loop
approximation) the following path integral expression for the
effective action ${\cal S}_{\rm {eff}}(\sigma,\pi)$, which indeed is a functional of the
bosonic $\sigma (t,\vec r)$ and $\pi_0 (t,\vec r)$ fields:
$$
\exp(\mathrm{i}{\cal S}_{\rm {eff}}(\sigma,\pi))=
  N'\int[d\bar q][dq]\exp\Bigl(\mathrm{i}\int\widetilde L\,dtd^2\vec r\Bigr),
$$
where
\begin{equation}
{\cal S}_{\rm {eff}}
(\sigma,\pi_0)
=-N\int dtd^2\vec r\left [\frac{\sigma^2(t,\vec r)+\pi_0^2(t,\vec r)}{4G}
\right ]+\widetilde {\cal S}_{\rm {eff}},
\label{F3}
\end{equation}
and $N'$ is a normalization constant. The fermion contribution to this effective action, i.e. the term
$\widetilde {\cal S}_{\rm {eff}}$ in (\ref{F3}), is given by:
\begin{equation}
\exp(\mathrm{i}\widetilde {\cal S}_{\rm {eff}})=N'\int [d\bar
q][dq]\exp\Bigl(\mathrm{i}\int\Big\{\bar q\big
[\gamma^\rho\mathrm{i}\partial_\rho +\mu\gamma^0+\sigma(t,\vec r) +\mathrm{i}\gamma^5\pi_0(t,\vec r)\big
]q\Big\}dtd^2\vec r\Bigr)=\det D.
 \label{F4}
\end{equation}
Here $D\,=\,\gamma^\rho\mathrm{i}\partial_\rho
+\mu\gamma^0+\sigma (t,\vec r)+i\gamma^5\pi_0 (t,\vec r)$. Now, using general formula
$\det D=\exp\left(Tr_{xsf}\ln D\right)$, where $Tr$ operation is supposed over coordinates (x),
spinor (s) and $O(N)$, i.e. flavor (f), indices, one can obtain the following expression for the effective action in the large-$N$ limit (mean-field approximation),
\begin{equation}
{\cal S}_{\rm {eff}}(\sigma,\pi_0)
=-N\int dtd^2\vec r\left[\frac{\sigma^2(t,\vec r)+\pi^2_0 (t,\vec r)}{4G}\right]-iN {\rm Tr}_{xs}\ln D.
\label{F6}
\end{equation}
The ground state expectation values  $\langle\sigma(t,\vec r)\rangle$ and $\langle\pi_0(t,\vec r)\rangle$ of
the composite bosonic fields $\sigma$ and $\pi_0$ are the solutions of the following
saddle point equations,
\begin{eqnarray}
\frac{\delta {\cal S}_{\rm {eff}}(\sigma,\pi_0)}{\delta\sigma (t,\vec r)}=0,~~~~~
\frac{\delta {\cal S}_{\rm {eff}}(\sigma,\pi_0)}{\delta\pi_0 (t,\vec r)}=0.
\label{F7}
\end{eqnarray}
We suppose that the system defined by the Lagrangian (\ref{F1}) is in the state with thermodynamic equilibrium, so the quantities
$\vev{\sigma(t,\vec r)}$ and $\vev{\pi_0(t,\vec r)}$ should not depend on time coordinate.
In vacuum, i.e. in the state corresponding to an empty space with zero particle density and zero
value of the chemical potential $\mu$, these quantities do not depend on spatial coordinates as
well, i.e. scalar and pseudoscalar fermion-antifermion condensates are homogeneous. But
even when studying the phase structure of dense quark matter, i.e., with non-zero $\mu$, it is
usually assumed that these condensates (order parameters) are still independent of the spatial
coordinates.

\subsection{Homogeneous ansatz}

So, first of all, let us also for simplicity suppose  that arguments of the effective
action functional  ${\cal S}_{\rm {eff}}(\sigma,\pi_0)$ (\ref{F6}) do not depend on spacetime
coordinates, i.e. $\sigma(t,\vec r)\equiv\sigma=const$ and $\pi_0(t,\vec r)\equiv\pi_0=const$.
Moreover, in this case in the Dirac operator $D$ of Eq. (\ref{F6}) there are no coordinate
dependency in all nonderivative terms, and we denote it as $D_0$, i.e. its coordinate matrix
element looks like $$D_0(x-y)=\left(\gamma^\rho\mathrm{i}\partial_\rho
+\mu\gamma^0+\sigma +i\gamma^5\pi_0\right)\delta (x-y)=
\int\frac{d^3p}{(2\pi)^3}\left(\not\!{\tilde p}+\sigma+i\gamma^5\pi_0\right)e^{-ip(x-y)}.$$
Hence, Fourier transform $\overline{D_0}(p)$ of $D_0(x-y)$ has the form $\overline{D_0}(p)=
\left(\not\!{\tilde p}+\sigma+i\gamma^5\pi_0\right)$. So
\begin{eqnarray}
S_0(x-y)\equiv D_0^{-1}=\left(\gamma^\rho\mathrm{i}\partial_\rho
+\mu\gamma^0+\sigma +i\gamma^5\pi_0\right)^{-1}
=\int\frac{d^3p}{(2\pi)^3}\frac{\tilde p_\nu\gamma^\nu-\sigma+i\gamma^5\pi_0}
{\tilde p^2-\sigma^2-\pi^2_0}e^{-ip(x-y)},\label{F8}
\end{eqnarray}
where $\tilde p=(p_0+\mu,p_1,p_2)$. Then TDP
$\Omega_\mu (\sigma,\pi_0)$ of the model in the homogeneous case is defined by the relation
\begin{eqnarray}
{\cal S}_{\rm {eff}}(\sigma,\pi_0)=-N\Omega_\mu (\sigma,\pi_0)\int d^3x
=-N\left[\frac{\sigma^2+\pi^2_0}{4G}\right]\int d^3x-iN {\rm Tr}_{xs}\ln D_0\nonumber\\=
-N\left\{\frac{\sigma^2+\pi^2_0}{4G}+i\int\frac{d^3p}{(2\pi)^3}
\ln(\det\overline{D_0}(p))\right\}\int d^3x.
\label{F9}
\end{eqnarray}
Here we use the notation $d^3x=dtd^2\vec r$ and take into account the relation (\ref{G7}) from
Appendix \ref{ApG}. $\blacktriangleright$ Note that $\overline{D_0}(p)=\left(\not\!{\tilde p}+\sigma+i\gamma^5\pi_0\right)$,
i.e. it is a 4$\times$4 matrix in spinor space. To find a determinant of this matrix, we use the following
relation $\det\overline{D_0}(p)=\det\left(\gamma^3\overline{D_0}(p)\gamma^3\right)=\det
\left(-\not\!{\tilde p}+\sigma-i\gamma^5\pi_0\right)$, where matrix $\gamma^3$ and its properties
are presented in Appendix \ref{ApA}. Hence
\begin{eqnarray}
\det\overline{D_0}(p)=\Big\{\det\big[\left(-\not\!{\tilde p}+\sigma-i\gamma^5\pi_0\right)
\left(\not\!{\tilde p}+\sigma+i\gamma^5\pi_0\right)\big]\Big\}^{1/2}=
\Big\{\det\left[(\sigma^2+\pi_0^2-\tilde p^2)1\!\!1_s\right]\Big\}^{1/2},\label{F10}
\end{eqnarray}
where $1\!\!1_s$ is the unit matrix in the four dimensional spinor space. It is evident that $\det\left[(\sigma^2+\pi_0^2-\tilde p^2)1\!\!1_s\right]=(\sigma^2+\pi_0^2-\tilde p^2)^4$. As a result we see that $\det\overline{D_0}(p)=(\sigma^2+\pi_0^2-\tilde p^2)^2$. Moreover, it is clear from Eq. (\ref{F9}) that the TDP
$\Omega_\mu (\sigma,\pi_0)$ looks like
\begin{eqnarray}
\Omega_\mu (\sigma,\pi_0)=\frac{\sigma^2+\pi_0^2}{4G}+i\int\frac{d^3p}{(2\pi)^3}\mathrm{ln}\Big\{\big[(p_0+\mu)^2-E^2\big]^2\Big\},
\label{F11}
\end{eqnarray}
where $E=\sqrt{\sigma^2+\pi_0^2+\vec p^2}$. $\blacktriangleleft$ Of course, the TDP (\ref{F11})
is not renormalized and depends effectively on a single field combination $M\equiv
\sqrt{\sigma^2+\pi_0^2}$, i.e. $\Omega_\mu (\sigma,\pi_0)\equiv\Omega_\mu (M)$. To get a finite
(renormalized) expression for it, it is useful to see that this unrenormalized TDP coincides with TDP (\ref{12}) $\Omega_\mu (M,b)$ at $b=0$. Its renormalized expression obtained in the framework of momentum cutoff regularization scheme is presented by Eq. (\ref{21}). There, instead of a bare coupling constant $G$, there appears a finite parameter $g$ with the dimension of the inverse mass. These quantities are connected by the relation (\ref{16}), $\frac 1{4G}\equiv \frac
1{4G(\Lambda)}=\frac{2\Lambda\ln(1+\sqrt{2})}{\pi^2}+\frac{1}{2g},$ where $\Lambda$ is the regularization cutoff parameter. Then, as it was shown in Ref. \cite{klimenko}, at $g>0$ chiral symmetry remains intact, but at $g<0$ it can be broken spontaneously. The saddle point (or gap) equations (\ref{F7}) take the form
\begin{eqnarray}
\frac{\partial \Omega_\mu (\sigma,\pi_0)}{\partial\sigma}\equiv \frac{\sigma}{2G}-
4\sigma i\int\frac{d^3p}{(2\pi)^3}\frac{1}{[(p_0+\mu)^2-E^2]}=0,\nonumber\\
\frac{\partial \Omega_\mu (\sigma,\pi_0)}{\partial\pi_0}\equiv \frac{\pi_0}{2G}-
4\pi_0 i\int\frac{d^3p}{(2\pi)^3}\frac{1}{[(p_0+\mu)^2-E^2]}=0.
\label{F12}
\end{eqnarray}
The coordinates of the GMP of the TDP $\Omega_\mu(\sigma,\pi_0)$ are no more, but the solution of the gap system (\ref{F12}). For simplicity, we will use the same notation $(\sigma,\pi_0)$ for the coordinates of the GMP of the TDP, as for its arguments.
It is clear from Eqs. (\ref{F12}) that there are four types of GMPs of the TDP, i) $(\sigma\ne 0,\pi_0\ne 0)$,
ii) $(\sigma\ne 0,\pi_0=0)$, iii) $(\sigma=0,\pi_0\ne 0)$, and iv) $ (\sigma=0,\pi_0=0)$.
For each of the GMPs i)--iii) the quantity $M_0\equiv\sqrt{\sigma^2+\pi_0^2}$ is the same, i.e. $M_0=-\pi/g$. These GMPs are $U_A(1)$ equivalent and can appear only at $g<0$. The parameter $M_0$ in this case is connected with coupling constant $G$ by the relation
\begin{eqnarray}
\frac{1}{2G}=4 i\int\frac{d^3p}{(2\pi)^3}\frac{1}{[(p_0+\mu)^2-M_0^2-p_1^2-p_2^2]}.
\label{F13}
\end{eqnarray}
As it is shown in Ref. \cite{klimenko}, at $g<0$ and $\mu<M_0$ the GMP of the TDP (\ref{F11}) can be one of i)--iii) types, but at $\mu>M_0$ it is of the form iv) even at $g<0$ (recall that at $g>0$ the GMP is also has the form iv)).

\subsection{Inhomogeneous fluctuations of the GMPs of the form i)--iii)}

Let us suppose that TDP $\Omega_\mu (\sigma,\pi_0)$ (\ref{F11}) has a global minimum at
the point $(\sigma,\pi_0)$ with $M_0=\sqrt{\sigma^2+\pi_0^2}\ne 0$, i.e. it is one of the forms i)--iii). In this case the relation $\mu<M_0$ must be carried out obligatotorily. Usually, ground state of the system is identified just with this
GMP, because in this point system has a least energy. Of course, this point is stable under
{\it homogeneous} fluctuations of order parameters $\sigma$ and $\pi_0$, which are, by definition, the GMP of
the TDP. Now we want to study the stability of this state with respect to small
{\it inhomogeneous} fluctuations. To this end we suppose that arguments $\sigma (t,\vec r)$ and
$\pi_0(t,\vec r)$ of the effective action ${\cal S}_{\rm {eff}}(\sigma(t,\vec r),\pi_0(t,\vec r))$ (\ref{F6}) has the following form
$$\sigma(t,\vec r)\equiv\sigma(x)=\sigma+\epsilon (x),~\pi_0(t,\vec r)\equiv\pi_0(x)=\pi_0+\omega (x),$$
where $\epsilon (x)\equiv\epsilon (t,x^1,x^2)$ and $\omega (x)\equiv\omega (t,x^1,x^2)$ are
small fluctuation of the homogeneous condensates $\sigma$ and $\pi_0$, respectively.
(Here and below we use an abbreviation $x$ for a set of three Minkowski coordinates,
$x\equiv (t,x^1,x^2)$.) Then
\begin{eqnarray}
\ln D=\ln \big[(\gamma^\rho\mathrm{i}\partial_\rho
+\mu\gamma^0+\sigma+i\gamma^5\pi_0)+\epsilon+i\gamma^5\omega\big]=\ln S_0^{-1}+\ln
\big[1+S_0\cdot (\epsilon+i\gamma^5\omega)\big],
\label{F14}
\end{eqnarray}
where $S_0^{-1}\,=\,\gamma^\rho\mathrm{i}\partial_\rho
+\mu\gamma^0+\sigma+i\gamma^5\pi_0$. Recall that (here we use the notation $\tilde p\equiv (p_0+\mu,p_1,p_2)$)
\begin{eqnarray}
S_0(x-y)=\int\frac{d^3p}{(2\pi)^3}\frac{\tilde p_\nu\gamma^\nu-\sigma+i\gamma^5\pi_0}
{\tilde p^2-\sigma^2-\pi_0^2}e^{-ip(x-y)}\equiv \int\frac{d^3p}{(2\pi)^3}\overline{S_0}(p)
e^{-ip(x-y)}.
\label{F15}
\end{eqnarray}
Expanding the functional ${\cal S}_{\rm {eff}}(\sigma(t,\vec r),\pi_0(t,\vec r))$ (\ref{F6})  in the neighborhood of the GMP $(\sigma,\pi_0)$ in a Taylor series with respect to the small fluctuations $\epsilon$ and $\omega$, we have
\begin{eqnarray}
{\cal S}_{\rm {eff}}
(\sigma+\epsilon(x),\pi_0+\omega(x))=
-\frac{N}{4G}\int d^3x \Big[\big(\sigma+\epsilon(x)\big)^2+\big(\pi_0+\omega(x)\big)^2\Big]-iN {\rm Tr}_{xs}\Big\{\ln S_0^{-1}+
\ln(1+S_0\cdot \big(\epsilon+i\gamma^5\omega\big)\Big\}\nonumber\\=-\frac{N}{4G}\int d^3x \big(\sigma^2+\pi_0^2\big)-iN
{\rm Tr}_{xs}\Big\{\ln S_0^{-1}\Big\}-\frac{N}{4G}\int d^3x \big(2\sigma\epsilon(x)+2\pi_0\omega(x)\big)
-iN {\rm Tr}_{xs}\Big\{S_0\cdot \big(\epsilon+i\gamma^5\omega\big)\Big\}\nonumber\\-\frac{N}{4G}\int d^3x
\big(\epsilon^2(x)+\omega^2(x)\big)+\frac{iN}{2} {\rm Tr}_{xs}\Big\{S_0\cdot \big(\epsilon+
i\gamma^5\omega\big)\cdot S_0\cdot \big(\epsilon+i\gamma^5\omega\big)\Big\}+\cdots\equiv
{\cal S}^{(0)}_{\rm {eff}}+{\cal S}^{(1)}_{\rm {eff}}+{\cal S}^{(2)}_{\rm {eff}}+\cdots,~~~
\label{F16}
\end{eqnarray}
where ${\cal S}^{(0)}_{\rm {eff}}=-N\Omega_\mu (\sigma,\pi_0)\int d^3x$ and the TDP $\Omega_\mu
(\sigma,\pi_0)$ is presented in Eq. (\ref{F11}). Since
${\rm Tr}_{xs}\Big\{S_0\cdot \big(\epsilon+i\gamma^5\omega\big)\Big\}=\int d^3x\cdot {\rm tr}_s
\Big\{S_0(0) \big(\epsilon (x)+i\gamma^5\omega(x)\big)\Big\}~$ (see Eq. (\ref{G10}) in Appendix \ref{ApG}),
it is easy to see that
\begin{eqnarray}
{\cal S}^{(1)}_{\rm {eff}}&=&-N\int d^3x\Big\{\epsilon (x)\Big[\frac{\sigma}{2G}-4i\sigma
\int\frac{d^3p}{(2\pi)^3}\frac{1}{\tilde p^2-\sigma^2-\pi_0^2}\Big]+\omega(x)\Big[\frac{\pi_0}{2G}-4i\pi_0
\int\frac{d^3p}{(2\pi)^3}\frac{1}{\tilde p^2-\sigma^2-\pi_0^2}\Big]\Big\}\equiv 0.~~
\label{F17}
\end{eqnarray}
The last identity is trivial since $\sigma$ and $\pi_0$ are the solution of one of the gap Eqs. (\ref{F12}).

Finally, due to a general relation (see Eq. (\ref{G1}) in Appendix \ref{ApG})
$${\rm Tr}_{xs}\Big\{S_0\cdot \big(\epsilon+i\gamma^5\omega\big)\cdot S_0\cdot \big(\epsilon+i\gamma^5
\omega\big)\Big\}
=\int d^3x\int d^3y~{\rm tr}_s\Big\{S_0(x-y)\big(\epsilon(y)+i\gamma^5\omega(y)\big)
S_0(y-x)\big(\epsilon(x)+i\gamma^5\omega(x)\big)\Big\},$$ we have for ${\cal S}^{(2)}_{\rm {eff}}$ term of Eq. (\ref{F16}) the following expression
\begin{eqnarray}
{\cal S}^{(2)}_{\rm {eff}}&=&-\frac{N}{4G}\int d^3x \big(\epsilon^2(x)+\omega^2(x)\big)+\frac{iN}{2}\int d^3x\int d^3y~\Big\{{\rm tr}_s\Big[S_0(x-y)S_0(y-x)\Big]\epsilon(y)\epsilon(x)\nonumber\\&+&
i\cdot {\rm tr}_s\Big[S_0(x-y)\gamma^5S_0(y-x)\Big]\omega(y)\epsilon(x)+i\cdot {\rm tr}_s\Big[S_0(x-y)S_0(y-x)\gamma^5\Big]\omega(x)\epsilon(y)\nonumber\\&-&{\rm tr}_s\Big[S_0(x-y)\gamma^5S_0(y-x)\gamma^5\Big]\omega(x)\omega(y)\Big\}.
\label{F18}
\end{eqnarray}

$\bullet$ Let us rewrite the last expression in terms of the Fourier transformations of the propagators (see in Eq. (\ref{F15})), as well as the fluctuation fields $\epsilon (x)$ and $\omega(x)$, i.e., we use in (\ref{F18}) the relation (\ref{F15}), as well as the following representation
\begin{eqnarray}
\epsilon (x)=\int\frac{d^3q_1}{(2\pi)^3}\overline\epsilon (q_1)e^{-iq_1x},~~~
\omega (y)=\int\frac{d^3q_2}{(2\pi)^3}\overline\omega (q_2)e^{-iq_2y}.
\label{F19}
\end{eqnarray}
In terms of Fourier transformations one of the $\gamma^5$ terms in Eq. (\ref{F18}) has the form
\begin{eqnarray}
T_{\gamma^51}\equiv\frac{iN}{2}\int d^3x\int d^3y~
i\cdot {\rm tr}_s\Big[S_0(x-y)\gamma^5S_0(y-x)\Big]\omega(y)\epsilon(x)=\frac{-N}{2}\int d^3x d^3y\int\frac{d^3p_1}{(2\pi)^3}\frac{d^3p_2}{(2\pi)^3}\frac{d^3q_1}{(2\pi)^3}\frac{d^3q_2}{(2\pi)^3}\times\nonumber\\{\rm tr}_s\Big\{\overline S(p_1)\gamma^5\overline
S(p_2)\Big\}\overline \epsilon (q_1)\overline \omega (q_2)e^{-ip_1(x-y)-ip_2(y-x)-iq_1x-iq_2y}.~~~~~~~~~~~~~~~~~~~~~~~~~~~~
\label{F20}
\end{eqnarray}
After integrations over $x$ and $y$, it takes the form
\begin{eqnarray}
T_{\gamma^51}&=&\frac{-N}{2}\int d^3q_1d^3q_2\frac{d^3p_1}{(2\pi)^3}\frac{d^3p_2}{(2\pi)^3}\delta^{(3)}(p_1-p_2+q_1)\delta^{(3)}(p_2-p_1+q_2)%\nonumber\\
{\rm tr}_s\Big\{\overline S(p_1)\gamma^5\overline
S(p_2)\Big\}\overline \epsilon (q_1)\overline \omega (q_2).
\label{F21}
\end{eqnarray}
Then, integrating in Eq. (\ref{F21}) over $q_1$ and $q_2$ with the help of $\delta$-functions, we have
\begin{eqnarray}
T_{\gamma^51}&=&\frac{-N}{2}\int \frac{d^3p_1}{(2\pi)^3}\frac{d^3p_2}{(2\pi)^3}{\rm tr}_s\Big\{\overline S(p_1)\gamma^5\overline
S(p_2)\Big\}\overline \epsilon (p_2-p_1)\overline \omega (p_1-p_2)\nonumber\\&=&\frac{-N}{2}\int \frac{d^3p}{(2\pi)^3}\frac{d^3q}{(2\pi)^3}{\rm tr}_s\Big\{\overline S(p)\gamma^5\overline
S(p+q)\Big\}\overline \epsilon (q)\overline \omega (-q)\nonumber\\&=&N4i\pi_0\sigma\int \frac{d^3p}{(2\pi)^3}\frac{d^3q}{(2\pi)^3}\frac{\overline \epsilon (q)\overline \omega (-q)}{[(p_0+\mu)^2-M_0^2-\vec p^2][(p_0+q_0+\mu)^2-M_0^2-(\vec p+\vec q)^2]}.
\label{F22}
\end{eqnarray}
Here we use the notation $M_0^2=\sigma^2+\pi_0^2$. We have already
mentioned above that in state of thermodynamic equilibrium the order parameters, i.e. the
quantities $\sigma (x),\pi_0(x)$ and $\epsilon (x),\omega(x)$, should not depend on time, i.e. on the coordinate $x^0$. As it is clear from Eq. (\ref{F19}), in terms of Fourier transforms $\overline \epsilon (q)$ and $\overline \omega (q)$ it means that
$\overline \epsilon (q)\equiv \overline \epsilon (q_0,q_1,q_2)=2\pi \delta (q_0)\overline \epsilon (q_1,q_2)$. Similarly, in this case we have $\overline \omega(q)\equiv \overline \omega (q_0,q_1,q_2)=2\pi \delta (q_0)\overline \omega (q_1,q_2)$. Moreover, since $\delta (q_0)=\int\frac{dx^0}{2\pi}e^{iq_0x^0}$, it is clear
that $\delta (q_0){\Big |}_{q_0=0}=\int\frac{dx^0}{2\pi} ~(*)$ (see also the remark at the end of Appendix \ref{ApG}). Supposing that in Eq. (\ref{F22})
\begin{eqnarray}
\overline \epsilon (q)\equiv \overline \epsilon (q_0,q_1,q_2)=2\pi \delta (q_0)\overline
\epsilon (q_1,q_2),~\overline \omega(-q)\equiv \overline \omega (-q_0,-q_1,-q_2)=2\pi \delta (-q_0)\overline \omega (-q_1,-q_2),
\label{0F23}
\end{eqnarray}
we have integrated there over $q_0$ and then taken into account the above relation $(*)$. As a result:
\begin{eqnarray}
T_{\gamma^51}=N4i\pi_0\sigma\int dx^0\int \frac{dq_1dq_2}{(2\pi)^2}\frac{d^3p}{(2\pi)^3}\frac{\overline \epsilon (q_1,q_2)\overline \omega (-q_1,-q_2)}{[(p_0+\mu)^2-M_0^2-\vec p^2][(p_0+\mu)^2-M_0^2-(\vec p+\vec q)^2]}.
\label{F024}
\end{eqnarray}

$\bullet$ The second $\gamma^5$ term in Eq. (\ref{F18}) has the form
\begin{eqnarray}
T_{\gamma^52}\equiv\frac{iN}{2}\int d^3x\int d^3y~
i\cdot {\rm tr}_s\Big[S_0(x-y)S_0(y-x)\gamma^5\Big]\omega(x)\epsilon(y).
\label{F23}
\end{eqnarray}
Using in Eq. (\ref{F23}) a well-known property of the  ${\rm tr}$ operation, i.e. ${\rm tr}
(A\cdot B\cdot C)={\rm tr}(C\cdot A\cdot B)$, and then carrying out a substitution of
integration variables there, i.e. $x\leftrightarrow y$, it is easy to see that
$T_{\gamma^52}=T_{\gamma^51}$. Changing variables in the integral (\ref{F024}), $\vec q\to -\vec q$, we get for $T_{\gamma^52}$ the expression
\begin{eqnarray}
T_{\gamma^52}=N4i\pi_0\sigma\int dx^0\int \frac{dq_1dq_2}{(2\pi)^2}\frac{d^3p}{(2\pi)^3}\frac{\overline \omega (q_1,q_2)\overline \epsilon (-q_1,-q_2)}{[(p_0+\mu)^2-M_0^2-\vec p^2][(p_0+\mu)^2-M_0^2-(\vec p-\vec q)^2]}.
\label{F023}
\end{eqnarray}

$\bullet$ The $T_{\epsilon\epsilon}$ term of Eq. (\ref{F18}) looks like (here $M_0^2=\sigma^2+\pi_0^2$)
\begin{eqnarray}
T_{\epsilon\epsilon}&=&-\frac{N}{4G}\int d^3x \epsilon^2(x)+\frac{iN}{2}\int d^3x\int d^3y~
\Big\{{\rm tr}_s\Big[S_0(x-y)S_0(y-x)\Big]\epsilon(y)\epsilon(x)\nonumber\\&=&
-N\int \frac{d^3q}{(2\pi)^3}
\overline \epsilon (q)\overline \epsilon (-q)\Big\{\frac{1}{4G}-\frac{i}{2}\int \frac{d^3p}{(2\pi)^3}{\rm tr}_s
\Big[\overline S(p)\overline S(p+q)\Big]\Big\}=-N\int \frac{d^3q}{(2\pi)^3}
\overline \epsilon (q)\overline \epsilon (-q)\Big\{\frac{1}{4G}\nonumber\\
&-&2i\int \frac{d^3p}{(2\pi)^3}\frac{(p_0+\mu)(p_0+q_0+\mu)-\vec p(\vec p+\vec q)+
\sigma^2-\pi_0^2}{[(p_0+\mu)^2-M_0^2-\vec p^2][(p_0+q_0+\mu)^2-M_0^2-(\vec p+\vec q)^2]}
\Big\}.
\label{F24}
\end{eqnarray}
If there is a state of thermodynamic equilibrium in the system, then $T_{\epsilon\epsilon}$
transforms into (below we use the notations $\overline\epsilon (\pm\vec q)\equiv\overline
\epsilon (\pm q_1,\pm q_2)$)
\begin{eqnarray}
T_{\epsilon\epsilon}&=&-N\int dx^0\int \frac{dq_1dq_2}{(2\pi)^2}
\overline \epsilon (\vec q)\overline \epsilon (-\vec q)\Big\{\frac{1}{4G}
-2i\int \frac{d^3p}{(2\pi)^3}\frac{(p_0+\mu)^2-\vec p(\vec p+\vec q)+
\sigma^2-\pi_0^2}{[(p_0+\mu)^2-M_0^2-\vec p^2][(p_0+\mu)^2-M_0^2-(\vec p+\vec q)^2]}
\Big\}\nonumber\\
&=&2iN\int dx^0\int \frac{dq_1dq_2}{(2\pi)^2}
\overline \epsilon (\vec q)\overline \epsilon (-\vec q)\int \frac{d^3p}{(2\pi)^3}\frac{2\sigma^2+
\vec q(\vec p+\vec q)}{[(p_0+\mu)^2-M^2-\vec p^2][(p_0+\mu)^2-M^2-(\vec p+\vec q)^2]}
\Big\}.
\label{F25}
\end{eqnarray}
The last expression in Eq. (\ref{F25}) is obtained after taking into account the equality
(\ref{F13}) for $G$. Finally, the $T_{\omega\omega}$ term of Eq. (\ref{F18}) is selected in the form (here $M_0^2=\sigma^2+\pi_0^2$)
\begin{eqnarray}
T_{\omega\omega}&=&-\frac{N}{4G}\int d^3x \omega^2(x)-\frac{iN}{2}\int d^3x\int d^3y~
\Big\{{\rm tr}_s\Big[S_0(x-y)\gamma^5S_0(y-x)\gamma^5\Big]\omega(y)\omega(x)\nonumber\\&=&
-N\int \frac{d^3q}{(2\pi)^3}
\overline \omega (q)\overline \omega (-q)\Big\{\frac{1}{4G}+\frac{i}{2}\int \frac{d^3p}{(2\pi)^3}{\rm tr}_s
\Big[\overline S(p)\gamma^5\overline S(p+q)\gamma^5\Big]\Big\}=-N\int \frac{d^3q}{(2\pi)^3}
\overline \omega (q)\overline \omega (-q)\Big\{\frac{1}{4G}\nonumber\\
&+&2i\int \frac{d^3p}{(2\pi)^3}\frac{-(p_0+\mu)(p_0+q_0+\mu)+\vec p(\vec p+\vec q)+
\sigma^2-\pi_0^2}{[(p_0+\mu)^2-M_0^2-\vec p^2][(p_0+q_0+\mu)^2-M_0^2-(\vec p+\vec q)^2]}
\Big\}.
\label{F26}
\end{eqnarray}
If a state of thermodynamic equilibrium is realized in a system, then conditions (\ref{0F23})
must be met. In this case (below we use the notations $\overline\omega (\pm\vec q)\equiv\overline
\omega (\pm q_1,\pm q_2)$)
\begin{eqnarray}
T_{\omega\omega}&=&-N\int dx^0\int \frac{dq_1dq_2}{(2\pi)^2}
\overline \omega (\vec q)\overline \omega (-\vec q)\Big\{\frac{1}{4G}
-2i\int \frac{d^3p}{(2\pi)^3}\frac{(p_0+\mu)^2-\vec p(\vec p+\vec q)-
\sigma^2+\pi_0^2}{[(p_0+\mu)^2-M^2-\vec p^2][(p_0+\mu)^2-M^2-(\vec p+\vec q)^2]}
\Big\}\nonumber\\&=&2iN\int dx^0\int \frac{dq_1dq_2}{(2\pi)^2}
\overline \omega (\vec q)\overline \omega (-\vec q)\int \frac{d^3p}{(2\pi)^3}\frac{2\pi_0^2+
\vec q(\vec p+\vec q)}{[(p_0+\mu)^2-M_0^2-\vec p^2][(p_0+\mu)^2-M_0^2-(\vec p+\vec q)^2]}
\Big\}.
\label{F27}
\end{eqnarray}
%This expression is very similar to Eq. (\ref{F25})

$\bullet$ Now, changing variables in the $p$ integrals (\ref{F024}), (\ref{F25}) and
(\ref{F27}), when $\vec p\to\vec p-\vec q/2$, and $\vec p\to \vec p+\vec q/2$ in Eq.
(\ref{F023}) and then ignoring the linear over $\vec p$ terms in the numerators of the obtained
$p$ integrals (see the corresponding remark just after Eq. (\ref{D26})), we have
\begin{eqnarray}
T_{\gamma^51}=2iN\int dx^0\int \frac{d^2\vec q}{(2\pi)^2}\overline \epsilon (-\vec q)
\overline \omega (\vec q)\int
\frac{d^3p}{(2\pi)^3}
\frac{2\sigma\pi_0}{[(p_0+\mu)^2-M_0^2-(\vec p-\vec q/2)^2][(p_0+\mu)^2-M_0^2-(\vec p+\vec q/2)^2]},\nonumber\\
T_{\gamma^52}=2iN\int dx^0\int \frac{d^2\vec q}{(2\pi)^2}\overline \epsilon (-\vec q)
\overline \omega (\vec q)\int
\frac{d^3p}{(2\pi)^3}
\frac{2\sigma\pi_0}{[(p_0+\mu)^2-M_0^2-(\vec p-\vec q/2)^2][(p_0+\mu)^2-M_0^2-(\vec p+\vec q/2)^2]},
\nonumber\\
T_{\epsilon\epsilon}
=2iN\int dx^0\int \frac{d^2\vec q}{(2\pi)^2}
\overline \epsilon (-\vec q)\overline \epsilon (\vec q)\int \frac{d^3p}{(2\pi)^3}
\frac{2\sigma^2+
\vec q^2/2}{[(p_0+\mu)^2-M_0^2-(\vec p-\vec q/2)^2][(p_0+\mu)^2-M_0^2-(\vec p+\vec q/2)^2]}
\Big\},\nonumber\\
T_{\omega\omega}=2iN\int dx^0\int \frac{d^2\vec q}{(2\pi)^2}
\overline \omega (-\vec q)\overline \omega (\vec q)\int \frac{d^3p}{(2\pi)^3}\frac{2\pi_0^2+
\vec q^2/2}{[(p_0+\mu)^2-M_0^2-(\vec p-\vec q/2)^2][(p_0+\mu)^2-M_0^2-(\vec p+\vec q/2)^2]}
\Big\}.
\label{F28}
\end{eqnarray}
Since ${\cal S}^{(2)}_{\rm {eff}}$ in Eq. (\ref{F18}) has the form ${\cal S}^{(2)}_{\rm {eff}}=T_{\gamma^51}+
T_{\gamma^52}+T_{\epsilon\epsilon}+T_{\omega\omega}$, we obtain, using Eqs. (\ref{F28}) for its components,
\begin{eqnarray}
{\cal S}^{(2)}_{\rm {eff}}(\epsilon,\omega)=2iN\int dx^0\int \frac{d^2\vec q}{(2\pi)^2}
\left (\begin{array}{c}
\overline \epsilon (-\vec q),
\overline \omega (-\vec q)
\end{array}\right )\left (\begin{array}{cc}
2\sigma^2+
\vec q^2/2,& 2\sigma\pi_0\\
2\sigma\pi_0,& 2\pi_0^2+\vec q^2/2
\end{array}\right )
\left (\begin{array}{cc}
\overline \epsilon (\vec q)\\
\overline \omega (\vec q)
\end{array}\right )F(\vec q),
\label{F29}
\end{eqnarray}
where
\begin{eqnarray}
F(\vec q)=\int \frac{d^3p}{(2\pi)^3}\frac{1}{[(p_0+\mu)^2-M_0^2-(\vec p-\vec q/2)^2][(p_0+\mu)^2-M_0^2-(\vec p+\vec q/2)^2]}.
\label{F30}
\end{eqnarray}
$\blacktriangleright$ It is evident that matrix $\Gamma\equiv\left (\begin{array}{cc}
2\sigma^2+
\vec q^2/2,& 2\sigma\pi_0\\
2\sigma\pi_0,& 2\pi_0^2+\vec q^2/2
\end{array}\right )$ has two positive eigenvalues $\lambda_1=\vec q^2/2$ and $\lambda_2=2\sigma^2+2\pi_0^2+ \vec q^2/2$. The corresponding eigenvectors are $\vec e_{\lambda_1}=\left(\begin{array}{cc}
-\pi_0\\ \sigma
\end{array}\right )$ and $\vec e_{\lambda_2}=\left(\begin{array}{cc}
\sigma\\ \pi_0
\end{array}\right )$, i.e. $\Gamma\cdot\vec e_{\lambda_1}=\lambda_1\vec e_{\lambda_1}$ and $\Gamma\cdot\vec e_{\lambda_2}=\lambda_2\vec e_{\lambda_2}$. Their scalar production is zero, i.e. $(\vec e_{\lambda_1},\vec e_{\lambda_2})=\vec e^T_{\lambda_1}\cdot \vec e_{\lambda_2}=0.$ In addition, we see that $(\vec e_{\lambda_1},\vec e_{\lambda_1})=(\vec e_{\lambda_2}\cdot \vec e_{\lambda_2})=M_0^2=\sigma^2+\pi_0^2.$ It is easy to show that
\begin{eqnarray}
\vec B(\vec q)\equiv \left (\begin{array}{cc}
\overline \epsilon (\vec q)\\
\overline \omega (\vec q)
\end{array}\right )=\vec e_{\lambda_1}\frac{-\overline \epsilon (\vec q)\pi_0+\overline \omega (\vec q)\sigma}{M_0^2}+\vec e_{\lambda_2}\frac{\overline \epsilon (\vec q)\sigma+\overline \omega (\vec q)\pi_0}{M_0^2}.\label{F31}
\end{eqnarray}
Moreover, in Eq. (\ref{F29}) there is a row-matrix $\left (\begin{array}{c}
\overline \epsilon (-\vec q),
\overline \omega (-\vec q)
\end{array}\right )$, which can be obtained from the vector $\vec B(\vec q)$ (\ref{F31})
by the action of the complex conjugation and transposition operations, i.e. $\left (\begin{array}{c}
\overline \epsilon (-\vec q),
\overline \omega (-\vec q)
\end{array}\right )=\vec B(\vec q)^{*T}=\left (\begin{array}{cc}
\overline \epsilon (\vec q)\\
\overline \omega (\vec q)
\end{array}\right )^{*T}$. $\blacktriangleleft$
As a result, if we do not take into account the function $F(\vec q)$ in Eq. (\ref{F29}),
then the remaining part of the integrand will take the form (evidently, it is a non-negative)
\begin{eqnarray}
\left (\begin{array}{c}
\overline \epsilon (-\vec q),
\overline \omega (-\vec q)
\end{array}\right )\left (\begin{array}{cc}
2\sigma^2+
\vec q^2/2,& 2\sigma\pi_0\\
2\sigma\pi_0,& 2\pi_0^2+\vec q^2/2
\end{array}\right )
\left (\begin{array}{cc}
\overline \epsilon (\vec q)\\
\overline \omega (\vec q)
\end{array}\right )=~~~~~~~~\nonumber\\
\vec B(\vec q)^{*T}\cdot\Gamma\cdot\vec B(\vec q)=
\lambda_1\frac{\big |-\overline \epsilon (\vec q)\pi_0+\overline \omega (\vec q)\sigma\big |^2}{M_0^2}+\lambda_2\frac{\big |\overline \epsilon (\vec q)\sigma+\overline \omega (\vec q)\pi_0\big |^2}{M_0^2}\ge 0.\label{F32}
\end{eqnarray}
Now the properties of the function $F(\vec q)$ (\ref{F30}) are in order. Since in our case $\mu<M_0$, the technique presented below Eq. (\ref{D24}) can be applied in order to integrate over $p_0$ in Eq. (\ref{F30}). As a result, we have
\begin{eqnarray}
{\cal S}^{(2)}_{\rm {eff}}(\epsilon,\omega)=-N\int dx^0\int \frac{d^2\vec q}{(2\pi)^2}\int \frac{d^2p}{(2\pi)^2}\frac{\lambda_1\big |\overline \epsilon (\vec q)\pi_0+\overline \omega (\vec q)\sigma\big |^2+\lambda_2\big |\overline \epsilon (\vec q)\sigma+\overline \omega (\vec q)\pi_0\big |^2}{M_0^2E_+E_-(E_++E_-)}.
\label{F33}
\end{eqnarray}
where $E_\pm=\sqrt{M_0^2+(\vec p\pm\vec q/2)^2}$.


\section{Task}
Now you can see that in section "Stability of homogeneous vacuum in 3D NJL-type model under small inhomogeneous
fluctuations" there has been done stability analysis of (2+1) NJL model.

Obtain the formula for $S^2$ in the similar way as it is in this section "Stability of homogeneous vacuum in 3D NJL-type model under small inhomogeneous
fluctuations" but in the case of the set up with magnetic field, discussed above with Lagrangian (1). The difference with section "Stability of homogeneous vacuum in 3D NJL-type model under small inhomogeneous
fluctuations" is that there is now homogeneous magnetic field B.


 the difference is mostly now that the expression for the propagator is different , it is in the non-zero constant magnetic field now. Find the expresion for the propagator in the setup of Lagrangian (1).

 First the calculations are the same as in the case of $B=0$.
Then substitute the propagator you will obtain into the expression for $S^2$ and you have to calculate traces over gamma matrices etc., try to simplify it (if possible perform integration ans sums etc).

Perform the calculation of $S^2$ in all situations, when M is generated then M is not gererated. At all values of $\mu$, not only $\mu$ is greater and smaller than M.

The expression for $S^{2}$ contains a term with \(G\), which can be eliminated from the gap equation if \(M\) is dynamically generated; however, if \(M\) is not generated and equals zero, renormalization is required—meaning that, if possible, one must isolate the divergent part and cancel it out with the \(G\)-term so that only \(g\) remains.

Then we have the $\sigma$ and $\pi$ and you may need to diagonalize the matrix 2*2 as it is done in section "Stability of homogeneous vacuum in 3D NJL-type model under small inhomogeneous
fluctuations"

For $S^2$ we would like to have maximally simplified expression, maybe expressed as integral or in some special functions. Then we want to explore the sign of $S^2$ in order to find if there is instability of homogeneous ground state with respect to inhomogeneous fluctuations.
\end{document}
